% ----- Consignes exo 1 ----- %
\begin{td-exo}[Complexité autour du \textsc{Vertex Cover}]\, % 1 
	On considère une variante du problème classique du \textsc{Vertex Cover}.
	Dans cette variante, intitulée \textsc{Connected Vertex Cover}, l'ensemble des sommets \(S\) choisis doit être connexe.
	
	Pour rappel, voici la définition des deux problèmes:
	
	\begin{problem}{\textsc{Vertex Cover} (\(VC\))}
    \Input{Un graphe \(G=(V,E)\) et un entier \(k\)}
    \Question{Existe-t-il un ensemble de sommets \(S \subseteq V\) de taille au plus \(k\) tel que chaque arête de \(E\) est adjacente à au moins un sommet de \(S\)?}
\end{problem}
	\begin{problem}{\textsc{Connected Vertex Cover} (\(CVC\))}
    \Input{Un graphe \(G=(V,E)\) et un entier \(k\)}
    \Question{Existe-t-il un ensemble de sommets \(S \subseteq V\) de taille au plus \(k\) tel que chaque arête de \(E\) est adjacente à au moins un sommet de \(S\) et que \(S\) est connexe?}
\end{problem}
	\begin{problem}{\textsc{Dominating Set} (\(DS\))}
    \Input{Un graphe \(G=(V,E)\) et un entier \(k\)}
    \Question{Existe-t-il un ensemble de sommets \(S \subseteq V\) de taille au plus \(k\) tel que chaque sommet de \(V\) est soit dans \(S\) soit adjacent à un sommet de \(S\)?}
\end{problem}

	\begin{enumerate}
		\item Montrer que \textsc{Connected Vertex Cover} est \(NP\)-complet.

		\item Montrer que \textsc{Vertex Cover} reste \(NP\)-complet pour un graphe de degré au plus \(3\).

		\textbf{Aide:} Transformer les sommets de degré \(4\).

		\item Montrer que \textsc{Dominating Set} est \(NP\)-complet.
		La preuve se fera à partir de \textsc{Vertex Cover} 

		\item Quelle est la complexité du problème \textsc{Vertex Cover} dans les graphes bipartis?

		\item Quelle est la complexité du problème \textsc{Vertex Cover} dans les graphes avec un degré maximum de \(2\)?
	\end{enumerate}
\end{td-exo}

% ----- Solutions exo 1 ----- %
\iftoggle{showsolutions}{ 
	\begin{td-sol}[]\ % 1
		
	\end{td-sol}
}{}
